\section{Subsystem visual evidence and independent software checks}
\label{app:des-visual-tests}
This appendix retains the offline checks and figures that extend the Design explanations. No training or gantry motion was performed during report revision. The original inputs, scripts and recorded outputs are preserved in the separate supporting-evidence package. The integrated qualification tests remain in Results.

\subsection{Actual capture variation and input resolution}
The three camera-roll inputs in Figure~\ref{fig:des-seg-variation} are \texttt{WIN\_20261028\_23\_23\_30\_Pro.jpg}, \texttt{WIN\_20261028\_23\_28\_34\_Pro.jpg} and \texttt{WIN\_20261028\_23\_32\_40\_Pro.jpg}. Each stored image is $3264\times2448$ pixels. The page visibly contains three handwritten text rows. Automatic axis decisions were retained, including the failure in the third capture. The observed crop counts were 4, 3 and 28, respectively. Shadows are visually different, but neither lux levels nor physical angles were measured. These observations therefore explain a design limitation rather than qualifying operation across a specified lighting range.

The native Owen01 page ($2448\times3264$) produced 20 crops for 20 labelled lines. A separate exploratory bench first fitted each of three original pages inside $1600\times1600$ pixels. From each reduced page it produced an unchanged view, a synthetic left-to-right multiplier from 0.42 to 1.00, a synthetic $+3^\circ$ rotation using bicubic interpolation on an expanded white canvas, and a quarter turn. Each was passed through the same automatic axis and segmentation functions. Table~\ref{tab:des-resample-cases} retains every case, including unsuccessful ones.

\begin{table}[!htbp]\centering\small
\begin{tabular}{lrrrrr}\hline
Page & Label lines & Reduced & Gradient & $+3^\circ$ & $90^\circ$ \\\hline
Owen & 20 & 19 & 20 & 11 & 19 \\
Dylan & 24 & 24 & 24 & 25 & 24 \\
Yeukita & 26 & 26 & 26 & 26 & 26 \\
\hline\end{tabular}
\caption{Exploratory crop counts after extra input reduction and digital perturbation. Equal counts do not establish correct ink ownership.}
\label{tab:des-resample-cases}
\end{table}

The reduced Owen case produced 19 crops and lost the usable first-line crop, whereas the native page produced 20. Its rotated reduced image produced only 11. The gradient image happened to restore the count to 20, which is not evidence that adding a shadow improves segmentation: thresholds, interpolation and component geometry changed together. These cases show why nominally scale-relative thresholds do not guarantee resolution invariance. A suitable follow-up would compare matched stroke ownership at controlled resolutions and improve capture guidance; the present bench does not supply pixel-level annotations or attribute every error to one operation.

\subsection{Independent component and seam references}
Algorithm~\ref{alg:des-reference-checks} states the two independent reference comparisons.
The run-based eight-connected component labeller was compared with an independently written flood-fill reference on 100 binary masks of size $9\times11$, using NumPy seed 2026 and foreground probabilities spanning 0.1--0.8. The test compared the full same-component relation between foreground pixel pairs, not arbitrary numerical label values. All 100 partitions agreed. Small random masks exercise diagonal contact and label equivalences, but do not demonstrate large-image speed or prove every possible topology.

For the seam solver, 20 random $4\times6$ cost arrays (seed 83) were checked against exhaustive enumeration of all row sequences with adjacent row changes of at most one. Both objectives included the implementation's $0.02|\Delta y|$ penalty. The maximum objective discrepancy was $0$. The exhaustive reference tests optimal cost on those arrays; it does not establish that the chosen image-derived cost always follows the correct handwritten line boundary.

\begin{designalgorithm}
\caption{Independent partition and seam verification}
\label{alg:des-reference-checks}
\prequire{Fixed random seeds, mask cases and small seam-cost arrays.}
\begin{designpseudo}
\pstate{0}{\textbf{for each} binary mask $M$ \textbf{do}}
\pstate{1}{$A\leftarrow\operatorname{LabelComponents}(M)$; $B\leftarrow\operatorname{FloodFill8}(M)$}
\pstate{1}{\textbf{assert} $[A_p=A_q]=[B_p=B_q]$ for every foreground pair $(p,q)$}
\pstate{0}{\textbf{end for}}
\pstate{0}{\textbf{for each} small cost array $C$ \textbf{do}}
\pstate{1}{$p^*\leftarrow\operatorname{LeastCostSeam}(C,1)$}
\pstate{1}{$\mathcal P\leftarrow$ all paths with adjacent row differences at most one}
\pstate{1}{$E(p)\leftarrow\sum_x C(p_x,x)+0.02\sum_x|p_x-p_{x-1}|$}
\pstate{1}{Record $|E(p^*)-\min_{p\in\mathcal P}E(p)|$}
\pstate{0}{\textbf{end for}}
\pstate{0}{\textbf{return} partition agreements and seam-cost discrepancies}
\end{designpseudo}
\end{designalgorithm}

\subsection{CTC probability and derivative check}
Algorithm~\ref{alg:des-ctc-reference} gives the exhaustive probability and derivative procedure.
The loss check used float64 logits of shape $5\times1\times4$, blank index 0 and target $[1,1]$. All 1024 frame paths were enumerated. Consecutive repeats were collapsed before blanks were removed, and only paths yielding the target contributed to its probability. This independently exercises the repeated-label restriction that a single uninterrupted run of label 1 cannot encode two labels. The path sum gave negative log likelihood 3.875974036731708; the dynamic programme gave 3.875974036731709. The absolute difference was $4.44\times10^{-16}$.

All 20 logits were perturbed individually with central-difference step $10^{-5}$. The maximum absolute discrepancy from \texttt{train.py::CtcLoss} was $3.63\times10^{-11}$. Logits were drawn after the fixed-seed component-mask draws in the bench; the script preserves that order. These checks concern the implemented CTC loss and derivative on a finite feasible target. They do not test impossible alignments, every transcript length, trained recognition accuracy or the full CNN--LSTM backward pass.

\begin{designalgorithm}
\caption{Exhaustive CTC likelihood and finite-difference check}
\label{alg:des-ctc-reference}
\prequire{Small logits $z$, target $\ell$, blank index and step $\varepsilon$.}
\begin{designpseudo}
\pstate{0}{$(L,g)\leftarrow\operatorname{CtcLoss}(z,\ell)$; $p\leftarrow\operatorname{Softmax}(z)$; $Z\leftarrow0$}
\pstate{0}{\textbf{for each} frame path $\pi\in\{0,\ldots,C-1\}^T$ \textbf{do}}
\pstate{1}{\textbf{if} $\mathcal B(\pi)=\ell$ \textbf{then} $Z\leftarrow Z+\prod_t p_t(\pi_t)$}
\pstate{0}{\textbf{end for}}
\pstate{0}{Record $|L+\log Z|$}
\pstate{0}{\textbf{for each} logit index $i$ \textbf{do}}
\pstate{1}{$g_i^{\rm ref}\leftarrow[L(z+\varepsilon e_i)-L(z-\varepsilon e_i)]/(2\varepsilon)$}
\pstate{1}{Record $|g_i-g_i^{\rm ref}|$}
\pstate{0}{\textbf{end for}}
\pstate{0}{\textbf{return} loss discrepancy and maximum derivative discrepancy}
\end{designpseudo}
\end{designalgorithm}

\subsection{Stored profiles, emitted files and physical page}
Figure~\ref{fig:des-style-files} uses the final Owen, Dylan, Yeukita and IAM 150 profiles with the same text, seed 2026, requested x-height 4\,mm and jitter 0.15. The left column shows generated pen paths. The right column reconstructs pen-down paths from each emitted G-code file with the independent parser; each panel is fitted separately for legibility. The normal converter was used, including its path ordering, joins and near-ring closure. All four generated jobs had coordinates within the configured bed bounds. The figures illustrate profile-dependent geometry and file conversion, rather than measuring similarity to a human writer.

\begin{figure}[!htbp]\centering
\includegraphics[width=0.98\linewidth]{Figures/design-style-gcode-examples.pdf}
\caption{Four profiles reconstructed from emitted G-code.}
\label{fig:des-style-files}
\end{figure}

Algorithm~\ref{alg:des-file-reference} describes the serialization check. Nine controlled fixtures combined scale factors 1, 3 and 30 with line offsets 0, 8.4 and 25.2\,mm. Ring closure and joining were disabled, minimum-segment filtering was set to zero, and the independently parsed paths were paired with \texttt{ToMachine}'s output. Each fixture retained path and vertex counts, with no clamped points. Maximum coordinate error was $0.0004$\,mm, within the $0.0005$\,mm per-coordinate rounding bound from three-decimal emission. This comparison checks serialization after placement; it does not independently prove that the placement equation itself is correct or that a motor reached a coordinate.

\begin{designalgorithm}
\caption{Controlled G-code serialization correspondence}
\label{alg:des-file-reference}
\prequire{Known multi-path trajectory and configurations with joins and closure disabled.}
\begin{designpseudo}
\pstate{0}{\textbf{for each} fixture configuration $c$ \textbf{do}}
\pstate{1}{$P\leftarrow\operatorname{ToMachine}(T,c)$}
\pstate{1}{$f\leftarrow\operatorname{WriteGcode}(T,c)$; $Q\leftarrow\operatorname{ParseGcode}(f)$}
\pstate{1}{Compare path counts and corresponding vertex counts}
\pstate{1}{Record $\max_i\|P_i-Q_i\|_\infty$ and clamped-point count}
\pstate{0}{\textbf{end for}}
\pstate{0}{\textbf{return} count agreements, rounding discrepancies and clipping records}
\end{designpseudo}
\end{designalgorithm}

Figure~\ref{fig:des-physical-page} is separately identified as a photograph of actual robot output, from the existing Owen reproduction sheet. The line crops are newly obtained from that stored photograph. This closes the visual explanation from a physical page to the segmentation interface, without presenting the earlier drawing as a newly executed hardware experiment. No exact same-job correspondence with Figure~\ref{fig:des-style-files} is claimed.

\begin{figure}[!htbp]\centering
\includegraphics[width=0.95\linewidth]{Figures/design-physical-page-example.pdf}
\caption{Photographed output and extracted line crops.}
\label{fig:des-physical-page}
\end{figure}

\subsection{Candidate-scoring configuration and scope}
The Owen trials use the joint reader and eleven-writer \texttt{pw170} checkpoint, with explicit 0.17 normalisation and uniform-ink rendering. They score character accuracy and requested-writer posterior before repair, independently of inherited server defaults. Exact inputs, readings and full-precision scores are retained in \texttt{Evidence/bench.json}.

\subsection{Post-design seam-coefficient sensitivity}
\label{app:des-seam-sweep}
A new one-at-a-time check examines the existing seam weights. Owen01 and Yeukita01 are reduced to a maximum dimension of 1600 pixels, then passed through orientation, conditioning and line-centre estimation. The first, middle and last gaps on each page are selected by index before inspecting outputs. Eight settings share these six gaps, step limit seven and ink weight four: midpoint weight varies at step penalty 0.02, then step penalty varies at midpoint weight 0.012. This is post-design sensitivity, not a joint search or recovered historical tuning.

The parameterised reference matches all six unmodified default paths. Table~\ref{tab:app-seam-sweep} reports raw-ink intersections, total row movement and mean midpoint distance in working pixels; the midpoint is a preference, not a labelled correct boundary.

\begin{table}[!htbp]\centering\small
\begin{tabular}{|c|c|r|r|r|}
\hline Midpoint weight&Step penalty&Ink intersections&Row movement&Mean offset / px\\\hline
0.000&0.020&2&35&4.0553\\\hline
0.006&0.020&2&38&0.5517\\\hline
0.012&0.020&2&42&0.3625\\\hline
0.024&0.020&2&44&0.3261\\\hline
0.048&0.020&2&47&0.2622\\\hline
0.012&0.000&2&48&0.2622\\\hline
0.012&0.010&2&44&0.3261\\\hline
0.012&0.040&2&38&0.5517\\\hline
\end{tabular}
\caption{New coefficient-sensitivity check across six fixed real-page gaps; the implemented setting is 0.012/0.020.}
\label{tab:app-seam-sweep}
\end{table}

Increasing midpoint weight from zero to 0.012 reduces mean offset from 4.055 to 0.363 pixels, while row movement rises from 35 to 42 pixels. Doubling the weight gives 0.326 pixels offset. Increasing step penalty has the opposite trade-off. Every setting retains two raw-mask and three closed-mask intersections. These cases therefore show competing effects without establishing a unique optimum, eliminating intersections or measuring recognition accuracy. Input/code hashes and all 48 runs are supplied; no pipeline coefficient was changed.
\FloatBarrier

\subsection{Real-input calculation traces}\label{app:des-real-calculations}
The numerical substitutions in Design were captured from offline processing of existing project inputs. They are execution traces of the stated calculations, not reconstructed historical training decisions or measurements of the mechanism. The complete settings, input/model hashes and full-precision values are stored in \texttt{Evidence/real\_calculations.json}; \texttt{trace\_real\_calculations.py} reproduces them without changing model weights or final profiles.

The page-conditioning trace uses Owen01 reduced to 1600 rows by 1200 columns, giving an 800-row by 600-column working image before deskew. The network and extraction trace instead uses the separately identified archived first original Owen01 crop ($156\times2170$), its known transcript and the joint reader. The latter produces $160\times1\times75$ log probabilities: the checkpoint head has 75 rows, consistent with 74 alphabet symbols and blank. The selected eleven-writer model uses explicit 0.17-width preparation for the two-crop vote. These inputs must not be mistaken for one continuous new whole-system test.

For the first convolution calculation, the actual input patch and learned kernel, listed in the same row-major order, are
\[
X=\begin{bmatrix}-1&-1&-1\\-0.552941&-0.380392&-0.568627\\0.866667&0.960784&0.866667\end{bmatrix},\qquad
W=\begin{bmatrix}0.213027&-0.215063&0.234474\\0.068398&0.083297&0.187499\\-0.312354&-0.291350&-0.221197\end{bmatrix}.
\]
Using the full-precision values with bias $-0.000435125$ gives $-1.15133214$ at output channel 0, row 34, column 462. Rounded values above explain the sum; the manifest retains the precise arrays needed to reproduce it. The subsequent batch-normalisation and LSTM substitutions follow this same crop through the reader.

The path and timing trace generates ``Writing with style.'' using the final Owen profile, seed 2026, requested x-height 4\,mm and jitter 0.15. The generated file is \texttt{Evidence/real\_trace\_owen.gcode}. Its default configuration supplies drawing/travel feeds 900/2400\,mm/min, nominal conversion 80 microsteps/mm and the 190.6 by 251.3\,mm workspace. No physical command was sent. The pulse-delay calculation uses one emitted segment at that nominal scale; actual homing counts remain a separate hardware measurement. The trace also records the pre-file coordinate mapping, distance accumulators, 56 pen transitions and the resulting 32.8\,s nominal budget.

The existing random-mask, exhaustive seam and small-lattice checks remain independent numerical tests. Their synthetic inputs are intentional test fixtures. They are separate from these real-input traces and from the qualification results.

\clearpage
\subsection{Visual checks of numerical training operations}\label{app:vt-gradient}
The retained visual-test record contains 249 scalar perturbation comparisons: 73 sampled convolution-weight entries, 73 bias comparisons, 73 input entries and all 30 logits of a small CTC lattice. The sampled entries need not be distinct parameters. Central differences use a step of $10^{-6}$ for the convolution and $10^{-5}$ for CTC, following the numerical-check principle in Equation~\ref{eq:des-gradcheck}. Figure~\ref{fig:vt-grad} compares the resulting numerical derivatives with the implemented backward operations. Relative difference is the absolute discrepancy divided by the sum of the two derivative magnitudes, with a $10^{-12}$ denominator floor.

The largest relative difference is $7.50\times10^{-7}$, among the sampled convolution inputs; the largest absolute input discrepancy is $7.26\times10^{-9}$. The CTC relative difference is at most $1.69\times10^{-9}$. These checks support the tested local derivatives under the fixed fixtures. They do not exercise every layer, parameter, floating-point regime or complete training run.

\begin{figure}[!htbp]\centering
\includegraphics[width=0.96\linewidth]{Figures/VisualTests/T38_gradient_check.pdf}
\caption{Analytic and finite-difference gradients.}
\label{fig:vt-grad}
\end{figure}

\subsection{Stroke-width quantisation in the validation sweep}\label{app:vt-width}
Figure~\ref{fig:vt-width} shows four real crops processed at the four sweep fractions. The image transformation depends on the resulting odd integer dilation kernel as well as the fraction. For Owen and Abhinav, 0.06 and 0.09 both select $k=3$ on the illustrated crop; for Dylan, 0.09 and 0.11 both select $k=5$. Identical kernels give identical normalised images for that same crop. This does not imply that two settings produce identical training datasets, since x-height varies across crops. The validation bars reproduce the supplied training sweep, not a new accuracy measurement. Its best observed setting remains 0.17, while a single run at each setting does not isolate width effects from training variability.

\begin{figure}[!htbp]\centering
\includegraphics[width=0.96\linewidth]{Figures/VisualTests/T40_stroke_width_sweep.pdf}
\caption{Stroke-width kernels and validation sweep.}
\label{fig:vt-width}
\end{figure}

\subsection{Glyph coordinates and measured connectedness}\label{app:vt-glyphgeometry}
Figure~\ref{fig:vt-coordinates} shows the actual Owen \texttt{V} used in the coordinate calculation in Design: left boundary 146 pixels, baseline 107 pixels, x-height 38 pixels and slant $-4.5^\circ$. The stored path uses the baseline as its origin, measures height upwards and expresses lengths in x-heights. The lead is stored separately from the path so that source-cell spacing is not confused with the first surviving ink point.

\begin{figure}[!htbp]\centering
\includegraphics[width=0.96\linewidth]{Figures/VisualTests/T43_glyph_coordinates.pdf}
\caption{Image and stored glyph coordinate frames.}
\label{fig:vt-coordinates}
\end{figure}

Figure~\ref{fig:vt-connectedness} makes the connectedness definition inspectable on the same line. Three of its 39 eligible adjacent non-space boundaries are crossed by a traced centreline, giving $3/39=0.076923$. This is the line-level example already used in Design. A crossing records connected source ink under the extraction convention; it does not reveal the writer's natural pen-lift sequence or establish a connection fraction for every page by itself.

\begin{figure}[!htbp]\centering
\includegraphics[width=0.96\linewidth]{Figures/VisualTests/T49_connectedness.pdf}
\caption{Observed connections between neighbouring characters.}
\label{fig:vt-connectedness}
\end{figure}

\FloatBarrier

\subsection{Extended figures from the Design section}
\label{app:extended-design-figures}
Figure~\ref{fig:vt-seamfield} extends Figure~\ref{fig:des-seam-example} with both selected gaps from Abhinav01. One was selected for the greatest ink saving and the other for the largest departure from the midpoint. These are demonstrations of the cost function, not average accuracy measurements.
\begin{figure}[!htbp]\centering
\includegraphics[width=0.96\linewidth]{Figures/VisualTests/T25_seam_cost_field.pdf}
\caption{Complete seam and cost-field examples.}
\label{fig:vt-seamfield}
\end{figure}

Figure~\ref{fig:vt-featurestack} extends the first-row example in Figure~\ref{fig:des-firstfeatures}. It retains the three selected channels at every stage and the final height reduction for the same Owen line. Independent display scaling makes local patterns visible without implying equal activation ranges between stages.
\begin{figure}[!htbp]\centering
\includegraphics[width=0.90\linewidth]{Figures/VisualTests/T32_feature_stack.pdf}
\caption{Complete convolutional feature sequence.}
\label{fig:vt-featurestack}
\end{figure}

